\section{Results}
\label{sec:clustering_agent_representation_results}

\input{figures/clustering_agent_representation/space_invaders_projection.tex}

Figure~\ref{fig:space_invaders_projection} previews the kind of explanation produced by \gls{car}: a two-dimensional projection of the \gls{ls} of a surrogate policy trained on Space Invaders, together with the resulting clusters and their corresponding observations.
Each cluster corresponds to a distinct, human-interpretable stage of the game, from ``few or no invaders have been neutralized'' to ``invaders descend, but their numbers are reduced.''
The remainder of this section quantitatively evaluates this behavior across several Atari benchmark environments \cite{Bellemare_2013}, which are widely used in \gls{rl} research \cite{mnih2013playingatarideepreinforcement, schulman2017proximalpolicyoptimizationalgorithms, zahavy2017graying}.
The environments considered in this study are:

\begin{itemize}
    \item \href{https://ale.farama.org/environments/breakout/}{\textbf{Breakout}}: the agent controls a paddle to bounce a ball and break as many bricks as possible, the ball speeding up after each bounce.
    \item \href{https://ale.farama.org/environments/mario_bros/}{\textbf{Mario Bros}}: the agent must defeat enemies, avoid fireballs, and collect mushrooms \cite{mario_bros_atari_2600}.
    \item \href{https://ale.farama.org/environments/pong/}{\textbf{Pong}}: a two-player game in which the agent tries to prevent the ball from passing its side while attempting to score against the opponent.
    \item \href{https://ale.farama.org/environments/space_invaders/}{\textbf{Space Invaders}}: the agent controls a cannon and must destroy waves of invaders before they reach the ground.
    \item \href{https://ale.farama.org/environments/surround/}{\textbf{Surround}}: a two-player grid-based game in which each move blocks the traversed cell, the goal being to trap the opponent.
\end{itemize}

In all these environments, the agent receives visual observations and operates in discrete action spaces.
The initial policies $\OriginalPolicy$ are trained using \gls{ppo} \cite{schulman2017proximalpolicyoptimizationalgorithms} with a deep actor-critic architecture.
The network consists of four convolutional layers with configurations $(16, 4, 2)$, $(32, 4, 2)$, $(64, 4, 2)$, $(128, 4, 2)$, where each tuple represents (number of channels, kernel size, stride), followed by three fully connected layers with 512, 256, and 128 units, respectively.
To ensure that these initial policies effectively learn to maximize rewards, we follow the recommendations of \cite{machado2017revisitingarcadelearningenvironment} and compare the performance of our initial policies against the benchmark scores reported in the same study to verify that they meet the required performance standards.

For each initial policy, three surrogate policies $\SurrogatePolicy{}$ sharing the same architecture are trained as described in Section~\ref{sec:clustering_agent_representation_method}.
These surrogate policies are randomly initialized and trained on approximately one million observation-action pairs each.

\begin{table}[t]
\centering
\resizebox{\linewidth}{!}{
    \begin{tblr}{
        colspec={|c|c|c|c|},
        hlines,
        vlines,
        row{1} = {font=\bfseries, halign=c, valign=m}
    }
        Environment     & {\textit{Behavior} \\ \textit{Reconstruction}} & {\textit{Abstraction} \\ \textit{Score}} & {\textit{Clustering} \\ \textit{Universality}} \\
        Breakout        & 0.83 $\pm$0.03 & 0.82 $\pm$0.07 & 0.48 $\pm$0.11 \\
        Mario Bros      & 0.99 $\pm$0.03 & 0.71 $\pm$0.14 & 0.55 $\pm$0.10 \\
        Ms.\ Pac-Man    & 1.00 $\pm$0.03 & 0.43 $\pm$0.09 & 0.65 $\pm$0.08 \\
        Pong            & 0.95 $\pm$0.01 & 0.99 $\pm$0.00 & 0.49 $\pm$0.05 \\
        Space Invaders  & 0.97 $\pm$0.07 & 0.69 $\pm$0.06 & 0.67 $\pm$0.07 \\
        Surround        & 1.00 $\pm$0.01 & 0.55 $\pm$0.07 & 0.61 $\pm$0.09 \\
    \end{tblr}
}
\caption{Evaluation of \BehaviorReconstructionName, \AbstractionScoreName and \ClusteringUniversalityName across the six Atari environments.}
\label{tab:environments_metrics}
\end{table}

\FloatBarrier

We first evaluate the ability of the surrogate policies to replicate the behavior of the initial policies (\BehaviorReconstructionName, Table~\ref{tab:environments_metrics}).
Across all environments, the surrogate policies achieve an average \BehaviorReconstructionName of $0.94 \pm 0.02$, indicating that they reliably reproduce the decisions of the initial policies over 100 evaluation episodes.

For the remainder of the analysis, a random subset of 50,000 observations is drawn from the training dataset and embedded into the \gls{ls} of each surrogate policy, using the representation extracted directly after the last convolutional layer.
This projection results in a point cloud $\textcolor{myblue}{\mathcal{E}}$.
For clustering, we opted for the simplest approach and applied \emphasize{K-means clustering} \cite{macqueen1967multivariate} to the point cloud, using the Euclidean distance between the projected points to form the clusters.

Across all environments, policies, and numbers of clusters, the average \AbstractionScoreName is $0.99 \pm 0.00$.
This indicates that the clusters formed in the \gls{ls} are not mere reflections of the raw pixel inputs or of the action distributions, and therefore capture a genuine abstraction of the agent's behavior, as defined in Definition~\ref{def:abstraction}.

\begin{figure}[t]
    \centering
    \resizebox{\textwidth}{!}{
    \begin{tikzpicture}
    \begin{axis}[
        xlabel={Number of clusters},
        ylabel={\ClusteringUniversalityName},
        grid=major,
        xmin=1, xmax=30,
        ymin=0, ymax=1,
        ytick distance=0.25,
        ymajorgrids=true,
        width=\textwidth,
        height=0.55\textwidth,
        set layers=standard,
        legend style={at={(0.5,-0.22)}, anchor=north, legend columns=3,
            /tikz/every even column/.append style={column sep=0.4cm}},
    ]
        \addplot[color=breakoutColor, mark=\breakoutMark, mark size=\breakoutMarkSize] coordinates {
            (2,0.99) (3,0.39) (4,0.37) (5,0.70) (6,0.43) (7,0.50) (8,0.44) (9,0.41) (10,0.45) (11,0.51) (12,0.41) (13,0.41) (14,0.49) (15,0.54) (16,0.46) (17,0.48) (18,0.46) (19,0.46) (20,0.50) (21,0.50) (22,0.46) (23,0.49) (24,0.47) (25,0.47) (26,0.46) (27,0.44) (28,0.44) (29,0.39) (30,0.46)
        };
        \addlegendentry{Breakout}
        \addplot[color=marioBrosColor, mark=\marioBrosMark, mark size=\marioBrosMarkSize] coordinates {
            (2,0.98) (3,0.80) (4,0.50) (5,0.62) (6,0.62) (7,0.59) (8,0.51) (9,0.57) (10,0.52) (11,0.53) (12,0.45) (13,0.51) (14,0.47) (15,0.49) (16,0.51) (17,0.51) (18,0.50) (19,0.49) (20,0.55) (21,0.53) (22,0.53) (23,0.53) (24,0.52) (25,0.52) (26,0.52) (27,0.54) (28,0.51) (29,0.51) (30,0.52)
        };
        \addlegendentry{Mario Bros}
        \addplot[color=msPacmanColor, mark=\msPacmanMark, mark size=\msPacmanMarkSize] coordinates {
            (2,0.91) (3,0.75) (4,0.62) (5,0.59) (6,0.67) (7,0.78) (8,0.77) (9,0.71) (10,0.79) (11,0.69) (12,0.66) (13,0.69) (14,0.62) (15,0.60) (16,0.59) (17,0.61) (18,0.60) (19,0.60) (20,0.64) (21,0.63) (22,0.60) (23,0.67) (24,0.56) (25,0.54) (26,0.60) (27,0.62) (28,0.55) (29,0.58) (30,0.64)
        };
        \addlegendentry{Ms.\ Pac-Man}
        \addplot[color=pongColor, mark=\pongMark, mark options={line width=1.5pt, scale=1}, mark size=\pongMarkSize] coordinates {
            (2,0.59) (3,0.59) (4,0.49) (5,0.46) (6,0.51) (7,0.57) (8,0.60) (9,0.54) (10,0.48) (11,0.50) (12,0.47) (13,0.48) (14,0.44) (15,0.47) (16,0.44) (17,0.46) (18,0.44) (19,0.47) (20,0.45) (21,0.45) (22,0.51) (23,0.49) (24,0.46) (25,0.49) (26,0.50) (27,0.49) (28,0.52) (29,0.42) (30,0.47)
        };
        \addlegendentry{Pong}
        \addplot[color=spaceInvadersColor, mark=\spaceInvadersMark, mark size=\spaceInvadersMarkSize] coordinates {
            (2,0.81) (3,0.72) (4,0.93) (5,0.69) (6,0.74) (7,0.76) (8,0.61) (9,0.62) (10,0.67) (11,0.59) (12,0.59) (13,0.56) (14,0.61) (15,0.63) (16,0.69) (17,0.64) (18,0.65) (19,0.67) (20,0.67) (21,0.66) (22,0.65) (23,0.64) (24,0.69) (25,0.67) (26,0.67) (27,0.70) (28,0.68) (29,0.68) (30,0.66)
        };
        \addlegendentry{Space Invaders}
        \addplot[color=surroundColor, mark=\surroundMark, mark size=\surroundMarkSize] coordinates {
            (2,0.85) (3,0.74) (4,0.80) (5,0.72) (6,0.67) (7,0.67) (8,0.67) (9,0.69) (10,0.66) (11,0.72) (12,0.59) (13,0.67) (14,0.64) (15,0.60) (16,0.54) (17,0.52) (18,0.55) (19,0.55) (20,0.53) (21,0.54) (22,0.55) (23,0.55) (24,0.54) (25,0.59) (26,0.54) (27,0.52) (28,0.53) (29,0.54) (30,0.55)
        };
        \addlegendentry{Surround}
        \addplot[domain=1:30, color=white, very thick, on layer=axis foreground]{0.5};
        \addplot[domain=1:30, samples=2, color=black, very thick, dashed, on layer=axis foreground]{0.5};
    \end{axis}
    \end{tikzpicture}
    }
    \caption{Evolution of \ClusteringUniversalityName with respect to the number of clusters.}
    \label{fig:clustering_universality_curves}
\end{figure}


Figure~\ref{fig:clustering_universality_curves} reports the evolution of \ClusteringUniversalityName with respect to the number of clusters.
From three clusters onward, all curves stabilize around $0.5$ in every environment.
Since \ClusteringUniversalityName is computed from the \gls{ari}, this result indicates a strong similarity in the clustering of the \gls{ls} across surrogate policies within each environment \cite{amezquita2020orchestrating}, regardless of the number of clusters chosen.

The combination of a high \AbstractionScoreName and a strong \ClusteringUniversalityName indicates that the latent clusters are both highly similar across independently trained surrogate policies and distinct from those obtained in the raw observation or action space.
These results support the \RepresentationConvergenceHypothesisName and confirm that the \gls{ls} captures high-level concepts, making it a reliable representation for analyzing agent behavior.
